% GENERATED FILE. Edit canonical lessons, then run npm run generate:books.
% canonical-source-hash: fnv1a64:657e307d715ffe4e
% canonical-lessons: SA-C112-padmam, SA-C112-naksatram, SA-C112-suryah, SA-C112-bhallukah, SA-C112-srgalah

\chapter{Lotus, Star, Sun, Bear, Jackal}
\label{ch:sa-lotus-star-sun-bear-jackal}
\hlchaptermodality{112}

\begin{chapteropening}
\textbf{By the end of this chapter:} \emph{I can say a lotus, a star, the sun, a bear and a jackal.}

Complete the last lesson of chapter 112: i can say a lotus, a star, the sun, a bear and a jackal.
\end{chapteropening}

\section[padmam]{\sk{पद्मम्} (padmam) — a lotus}

\label{lesson:SA-C112-padmam}



\begin{quote}
Before the new one: say the Sanskrit for a leaf, then the Sanskrit for a flower.
\end{quote}

\subsection*{You'll want to know: \sk{पद्मम्}}
\textbf{\sk{पद्मम्}} — \emph{padmam} — \textquotedblleft{}a lotus\textquotedblright{}.

\begin{grammarlens}[title={What you've built}]
One more word for this chapter.
\end{grammarlens}

\subsection*{Your turn}
\begin{itemize}
\raggedright
  \item \emph{Say it:} \emph{padmam}
  \item \emph{Say it:} \emph{padmam}, once more
  \item \emph{Say it:} say \emph{padmam}
  \item \emph{Recall:} say the Sanskrit for a leaf, then the Sanskrit for a flower, then say \emph{padmam} again
\end{itemize}

\begin{tcolorbox}[breakable,colback=teal!4,colframe=teal!35!black,title={Before you move on}]
What does \sk{पद्मम्} mean? (\textquotedblleft{}A lotus\textquotedblright{}.) Say it once more.
\end{tcolorbox}
\section[nakṣatram]{\sk{नक्षत्रम्} (nakṣatram) — a star}

\label{lesson:SA-C112-naksatram}



\begin{quote}
Before the new one: say the Sanskrit for a flower, then the Sanskrit for a lotus.
\end{quote}

\subsection*{You'll want to know: \sk{नक्षत्रम्}}
\textbf{\sk{नक्षत्रम्}} — \emph{nakṣatram} — \textquotedblleft{}a star\textquotedblright{}.

\begin{grammarlens}[title={What you've built}]
One more word for this chapter.
\end{grammarlens}

\subsection*{Your turn}
\begin{itemize}
\raggedright
  \item \emph{Say it:} \emph{nakṣatram}
  \item \emph{Say it:} \emph{nakṣatram}, once more
  \item \emph{Say it:} say \emph{nakṣatram}
  \item \emph{Recall:} say the Sanskrit for a flower, then the Sanskrit for a lotus, then say \emph{nakṣatram} again
\end{itemize}

\begin{tcolorbox}[breakable,colback=teal!4,colframe=teal!35!black,title={Before you move on}]
What does \sk{नक्षत्रम्} mean? (\textquotedblleft{}A star\textquotedblright{}.) Say it once more.
\end{tcolorbox}
\section[sūryaḥ]{\sk{सूर्यः} (sūryaḥ) — the sun}

\label{lesson:SA-C112-suryah}



\begin{quote}
Before the new one: say the Sanskrit for a lotus, then the Sanskrit for a star.
\end{quote}

\subsection*{You'll want to know: \sk{सूर्यः}}
\textbf{\sk{सूर्यः}} — \emph{sūryaḥ} — \textquotedblleft{}the sun\textquotedblright{}.

\begin{grammarlens}[title={What you've built}]
One more word for this chapter.
\end{grammarlens}

\subsection*{Your turn}
\begin{itemize}
\raggedright
  \item \emph{Say it:} \emph{sūryaḥ}
  \item \emph{Say it:} \emph{sūryaḥ}, once more
  \item \emph{Say it:} say \emph{sūryaḥ}
  \item \emph{Recall:} say the Sanskrit for a lotus, then the Sanskrit for a star, then say \emph{sūryaḥ} again
\end{itemize}

\begin{tcolorbox}[breakable,colback=teal!4,colframe=teal!35!black,title={Before you move on}]
What does \sk{सूर्यः} mean? (\textquotedblleft{}The sun\textquotedblright{}.) Say it once more.
\end{tcolorbox}
\section[bhallūkaḥ]{\sk{भल्लूकः} (bhallūkaḥ) — a bear}

\label{lesson:SA-C112-bhallukah}



\begin{quote}
Before the new one: say the Sanskrit for a star, then the Sanskrit for the sun.
\end{quote}

\subsection*{You'll want to know: \sk{भल्लूकः}}
\textbf{\sk{भल्लूकः}} — \emph{bhallūkaḥ} — \textquotedblleft{}a bear\textquotedblright{}.

\begin{grammarlens}[title={What you've built}]
One more word for this chapter.
\end{grammarlens}

\subsection*{Your turn}
\begin{itemize}
\raggedright
  \item \emph{Say it:} \emph{bhallūkaḥ}
  \item \emph{Say it:} \emph{bhallūkaḥ}, once more
  \item \emph{Say it:} say \emph{bhallūkaḥ}
  \item \emph{Recall:} say the Sanskrit for a star, then the Sanskrit for the sun, then say \emph{bhallūkaḥ} again
\end{itemize}

\begin{tcolorbox}[breakable,colback=teal!4,colframe=teal!35!black,title={Before you move on}]
What does \sk{भल्लूकः} mean? (\textquotedblleft{}A bear\textquotedblright{}.) Say it once more.
\end{tcolorbox}
\section[śṛgālaḥ]{\sk{शृगालः} (śṛgālaḥ) — a jackal}

\label{lesson:SA-C112-srgalah}



\begin{quote}
Before the new one: say the Sanskrit for the sun, then the Sanskrit for a bear.
\end{quote}

\subsection*{You'll want to know: \sk{शृगालः}}
\textbf{\sk{शृगालः}} — \emph{śṛgālaḥ} — \textquotedblleft{}a jackal\textquotedblright{}.

\begin{grammarlens}[title={What you've built}]
One more word for this chapter.
\end{grammarlens}

\subsection*{Your turn}
\begin{itemize}
\raggedright
  \item \emph{Say it:} \emph{śṛgālaḥ}
  \item \emph{Say it:} \emph{śṛgālaḥ}, once more
  \item \emph{Say it:} say \emph{śṛgālaḥ}
  \item \emph{Recall:} say the Sanskrit for the sun, then the Sanskrit for a bear, then say \emph{śṛgālaḥ} again
\end{itemize}

\begin{tcolorbox}[breakable,colback=teal!4,colframe=teal!35!black,title={Before you move on}]
What does \sk{शृगालः} mean? (\textquotedblleft{}A jackal\textquotedblright{}.) Say it once more.
\end{tcolorbox}
